\subsection*{Section 7: Improved Upper Bound via Small-Clique Graph Coloring}

The following theorem improves the standard upper bound $M(n)\le \tau(n)+1$ by a polynomial factor, where $\tau(n)$ denotes the $n$-dimensional kissing number.

\begin{theorem}[Improved Upper Bound]\label{thm:upper-bound}
There exists an absolute constant $C>0$ such that, for all sufficiently large $n$,
\[
  M(n) \;\le\; C\,\tau(n)\,\sqrt{\frac{\ln(n+2)}{\ln \tau(n)}}.
\]
In particular, using the Kabatiansky--Levenshtein bound $\tau(n)\le 2^{0.401n(1+o(1))}$,
\[
  M(n) \;=\; O\!\left(\tau(n)\cdot\sqrt{\frac{\log n}{n}}\right)
  \;=\; 2^{0.401n(1+o(1))}\cdot O\!\left(\sqrt{\frac{\log n}{n}}\right).
\]
This is a polynomial-factor improvement over the greedy bound $M(n)\le\tau(n)+1$.
\end{theorem}

The proof rests on three ingredients, which we state as lemmas.

\begin{lemma}[de Bruijn--Erd\H{o}s Compactness]\label{lem:compactness}
$M(n) = \sup\bigl\{\chi(G_X) : X\subset\mathbb{R}^n\text{ finite},\;\|x-y\|\ge 1\ \forall x\ne y\in X\bigr\}.$
In other words, $M(n)$ is achieved as a supremum over \emph{finite} packings.
\end{lemma}

\begin{proof}
Let $\mathcal{P}$ be any unit sphere packing in $\mathbb{R}^n$ (possibly infinite), and let $G=G_{\mathcal{P}}$ be its tangency graph.  Every finite subgraph of $G$ is the tangency graph of a finite sub-packing, hence has chromatic number at most $\sup_{\text{finite }X}\chi(G_X)=:K$.  By the de Bruijn--Erd\H{o}s theorem~\cite{dBE1951}, an infinite graph whose every finite subgraph is $K$-colorable is itself $K$-colorable, so $\chi(G)\le K$.  Hence $M(n)\le K$; the reverse inequality $M(n)\ge K$ is trivial.
\end{proof}

\begin{lemma}[Geometric Graph Bounds]\label{lem:graph-bounds}
For every finite unit sphere packing $X\subset\mathbb{R}^n$ (centres with pairwise distance $\ge 1$), the tangency graph $G_X$ satisfies:
\begin{enumerate}
  \item[\textup{(a)}] \textup{(Degree bound)} $\Delta(G_X)\le\tau(n)$, the $n$-dimensional kissing number.
  \item[\textup{(b)}] \textup{(Clique bound)} $\omega(G_X)\le n+1$.
\end{enumerate}
\end{lemma}

\begin{proof}
\textbf{(a).} The neighbours of a centre $v\in X$ in $G_X$ are the centres $u\in X$ with $\|u-v\|=1$, all lying on the unit sphere $S^{n-1}(v)$ of radius $1$ centred at $v$, and pairwise at distance $\ge 1$.  Translating to $v=0$, this is a set of unit-sphere points with pairwise angular separation $\ge 60°$, i.e.\ a kissing configuration.  By definition of $\tau(n)$, the size of any kissing configuration is at most $\tau(n)$.

\textbf{(b).} A $k$-clique in $G_X$ is a set $\{x_1,\ldots,x_k\}\subset X$ with $\|x_i-x_j\|=1$ for all $i\ne j$ — an equidistant set of diameter $1$ in $\mathbb{R}^n$.  Setting $y_i=x_i-x_1$ for $i\ge 2$, we have $\|y_i\|=1$ and
\[
  \|y_i-y_j\|^2 = \|x_i-x_j\|^2 = 1, \quad i\ne j\ge 2,
\]
which gives $\langle y_i,y_j\rangle = \tfrac{1}{2}$ for all $i\ne j\ge 2$.  The Gram matrix $M$ of $y_2,\ldots,y_k$ satisfies $M_{ii}=1$ and $M_{ij}=\tfrac{1}{2}$ for $i\ne j$, i.e.\ $M=\tfrac{1}{2}I+\tfrac{1}{2}J$ where $J$ is the $(k-1)\times(k-1)$ all-ones matrix.  The eigenvalues of $M$ are $\tfrac{k}{2}$ (once) and $\tfrac{1}{2}$ ($k-2$ times), all positive, so $M$ is positive definite and has rank $k-1$.  Since $y_2,\ldots,y_k\in\mathbb{R}^n$, we need $k-1\le n$, i.e.\ $k\le n+1$.
\end{proof}

\begin{lemma}[Bonamy--Kelly--Nelson--Postle, 2022~{\cite{BKNP2022}}]\label{lem:bknp}
There exists an absolute constant $C_0>0$ such that the following holds.  For all $\omega\ge 3$ and all $\Delta\ge \omega^{C_0}$, every graph $G$ with $\Delta(G)\le\Delta$ and $\omega(G)\le\omega$ satisfies
\[
  \chi(G)\;\le\; 72\,\Delta\,\sqrt{\frac{\ln\omega}{\ln\Delta}}.
\]
\end{lemma}

\begin{proof}[Proof of Theorem~\ref{thm:upper-bound}]
By Lemma~\ref{lem:compactness}, $M(n)=\sup_{\text{finite }X}\chi(G_X)$.  Fix a finite unit sphere packing $X\subset\mathbb{R}^n$.  By Lemma~\ref{lem:graph-bounds}, $G_X$ has $\Delta(G_X)\le\tau(n)$ and $\omega(G_X)\le n+1$.

We verify the hypothesis of Lemma~\ref{lem:bknp} with $\Delta=\tau(n)$ and $\omega=n+2$: since $\tau(n)\ge 2^{0.2075n(1+o(1))}$ by the Kabatiansky--Levenshtein lower bound~\cite{KL1978}, we have $\ln\tau(n)\ge 0.2075n\ln 2$, while $\ln(n+2)^{C_0}=O(n)$, so $\tau(n)\ge(n+2)^{C_0}$ for all sufficiently large $n$.  Hence
\[
  \chi(G_X) \;\le\; 72\,\tau(n)\,\sqrt{\frac{\ln(n+2)}{\ln\tau(n)}}.
\]
Taking the supremum over all finite $X$ gives the first bound.

For the asymptotic form, since $\ln\tau(n)\le 0.401n\ln 2+o(n)$ (upper KL bound~\cite{KL1978}),
\[
  \sqrt{\frac{\ln(n+2)}{\ln\tau(n)}} = O\!\left(\sqrt{\frac{\log n}{n}}\right),
\]
giving $M(n)=O(\tau(n)\sqrt{\log n/n})=2^{0.401n(1+o(1))}\cdot O(\sqrt{\log n/n})$.

To see this is an improvement over greedy: the ratio of the new bound to $\tau(n)$ is $O(\sqrt{\log n/n})\to 0$, so the new bound is $o(\tau(n))$ and hence $o(M_{\mathrm{greedy}}(n))$.
\end{proof}

\begin{remark}
The key geometric input is $\omega(G_X)\le n+1$, which contrasts sharply with $\tau(n)\sim 2^{0.401n}$.  The logarithmic ratio $\ln\omega/\ln\Delta\sim (\log n)/(0.401n)\to 0$ is what drives the improvement.  The de Bruijn--Erd\H{o}s step (Lemma~\ref{lem:compactness}) is essential: the Bonamy--Kelly--Nelson--Postle theorem is a finite-graph result, and the transfer to arbitrary (possibly infinite) packings requires the compactness principle.
\end{remark}

\subsection*{Section 7b: Reed's Bound — A Sharper Estimate for Moderate Dimensions}

The BKNP bound in Theorem~\ref{thm:upper-bound} is asymptotically the strongest known result, but for moderate dimensions its large absolute constant ($C_0 = 72$) makes it weaker than a classical alternative.  We record this alternative, which gives a cleaner closed-form bound superior to BKNP in every dimension $n$ arising in practice.

\begin{theorem}[Reed's Bound for Sphere Packings]\label{thm:reed}
For all sufficiently large $n$,
\[
  M(n) \;\le\; \left\lceil\frac{\tau(n)+n+2}{2}\right\rceil.
\]
\end{theorem}

\begin{proof}
By Lemma~\ref{lem:compactness}, it suffices to bound $\chi(G_X)$ for every finite packing tangency graph $G_X$.  By Lemma~\ref{lem:graph-bounds}, $G_X$ has $\Delta(G_X)\le\tau(n)$ and $\omega(G_X)\le n+1$.  Reed's theorem~\cite{Reed1998} states that for every graph $G$ with $\Delta(G)\ge\Delta_0(\omega(G))$ (a threshold depending on the clique number),
\[
  \chi(G)\;\le\;\left\lceil\frac{\Delta(G)+\omega(G)+1}{2}\right\rceil.
\]
Since $\tau(n)\ge 2^{0.2075n}$ grows exponentially while $\omega(G_X)\le n+1$ is linear, the condition $\Delta\ge\Delta_0(\omega)$ holds for all sufficiently large $n$.  Substituting $\Delta\le\tau(n)$ and $\omega\le n+1$ gives
\[
  \chi(G_X)\;\le\;\left\lceil\frac{\tau(n)+n+2}{2}\right\rceil.
\]
Taking the supremum over all finite $X$ yields the result.
\end{proof}

\begin{remark}[Comparison of the Two Upper Bounds]\label{rem:comparison}
Theorem~\ref{thm:upper-bound} (BKNP) gives $M(n)\le C\tau(n)\sqrt{\log n/n}$, while Theorem~\ref{thm:reed} gives $M(n)\le\tau(n)/2 + O(n)$.  For large $n$:
\[
  \text{BKNP} \;\sim\; C\tau(n)\sqrt{\frac{\log n}{n}} \;=\; o(\tau(n)),
  \qquad
  \text{Reed} \;\sim\; \frac{\tau(n)}{2}.
\]
BKNP is asymptotically stronger ($o(\tau(n))$ vs.\ $\tau(n)/2$), but the constant $C=72$ means Reed's bound dominates in every dimension accessible by computation.  Concretely, for $n\le 10^5$ and $\tau(n)\le 2^{0.401n}$, Reed gives $M(n)\le\tau(n)/2+O(n)$, whereas BKNP gives $72\tau(n)\sqrt{\log n/n}>\tau(n)/2$ (since $72^2\log n/n > 1/4$ for $n\le 10^5$).  The crossover where BKNP beats Reed occurs around $n\approx 4\ln(144^2)\cdot n/\ln n$, i.e.\ $n$ of order $10^6$ or larger.

The combined best-known upper bound is therefore:
\[
  M(n)\;\le\;\min\!\left(\left\lceil\tfrac{\tau(n)+n+2}{2}\right\rceil,\;C\tau(n)\sqrt{\tfrac{\log n}{n}}\right).
\]
For $n\le 10^5$ the Reed bound dominates; for $n\to\infty$ the BKNP bound dominates.

\textbf{Selected values} using exact kissing numbers where known and $\tau(n)\le 2^{0.401n}$ otherwise:

\smallskip
\begin{center}
\begin{tabular}{r|r|r|r|r}
$n$ & $\tau(n)$ & Greedy $\tau+1$ & Reed $\lceil(\tau+n+2)/2\rceil$ & Improvement factor\\
\hline
$3$ & $12$ & $13$ & $9$ & $1.44\times$\\
$4$ & $24$ & $25$ & $15$ & $1.67\times$\\
$8$ & $240$ & $241$ & $125$ & $1.93\times$\\
$24$ & $196560$ & $196561$ & $98293$ & $\approx 2\times$\\
$n\to\infty$ & $\tau(n)$ & $\tau(n)+1$ & $\tau(n)/2+O(n)$ & $\to 2\times$\\
\end{tabular}
\end{center}
\smallskip

Reed's bound thus cuts the greedy upper bound by a factor approaching $2$, independent of $n$.  The BKNP bound cuts it by a factor $\sqrt{n/\log n}\to\infty$, but only asymptotically.
\end{remark}

\begin{remark}[Barriers to Further Improvement]
Both bounds above ultimately rely on $\tau(n)$, the kissing number, as the main input.  Any improvement of the form $M(n)\le\tau(n)^{1-c}$ for a fixed $c>0$ would require structural properties of contact graphs beyond the degree and clique bounds used here.  Candidate approaches include:
\begin{itemize}
  \item \textbf{Lovász theta / SDP}: The Lov\'asz $\vartheta$-function of the contact graph is related to the Cohn--Elkies LP bound for sphere packings, and might yield $M(n)\le\tau(n)/n^{\varepsilon}$ for some $\varepsilon>0$.
  \item \textbf{Entropy compression in neighborhoods}: Since each neighborhood $N(v)$ is itself a kissing configuration (a unit sphere packing on $S^{n-1}$), the local chromatic structure is recursive and may admit a stronger coloring argument.
  \item \textbf{Improving $\tau(n)$}: Any improvement to the Kabatiansky--Levenshtein kissing number bound would immediately improve both $M(n)$ bounds above.  No improvement has been found since 1978.
\end{itemize}
For the \emph{lower} bound, the barrier theorem (Section~6) shows the existing spectral framework is stuck at surplus $O(n^{2/3})$; beating this requires a new geometric construction.  Consequently, the gap between the polynomial lower bound $M(n)\ge n+\Omega(n^{2/3})$ and the exponential upper bound $M(n)\le O(\tau(n))$ remains the central open problem.
\end{remark}
